% GENERATED FILE. Edit canonical lessons, then run npm run generate:books.
% canonical-source-hash: fnv1a64:c90d02ceec40d5cd
% canonical-lessons: AR-C517-dakhkhana, AR-C517-ishtaha, AR-C517-admana, AR-C517-jaa, AR-C517-shabia

\chapter{Habits good and bad}
\label{ch:ar-habits-good-and-bad}
\hlchaptermodality{517}

\begin{chapteropening}
\textbf{By the end of this chapter:} \emph{I can say that someone smokes, became addicted, got hungry, was full or craved something.}

Complete the last lesson of chapter 517: I can say that someone smokes, became addicted, got hungry, was full or craved something.
\end{chapteropening}

\section[\texorpdfstring{\ar{دخن} (dakhkhana)}{دخن (dakhkhana)}]{\ar{دخن} (dakhkhana) — to smoke tobacco}

\label{lesson:AR-C517-dakhkhana}



\begin{quote}
Before the new one: say the Arabic for to have a cough, then the Arabic for a sprain.
\end{quote}

\subsection*{You'll want to know: \ar{دخن}}
\textbf{\ar{دخن}} — \emph{dakhkhana} — \textquotedblleft{}to smoke tobacco\textquotedblright{}.

\begin{quote}
\ar{لا} \ar{يدخن} \ar{أبي} \ar{في} \ar{البيت}.
\end{quote}

\begin{quote}
\emph{lā yudakhkhinu abī fī al-bayt.}
\end{quote}

\begin{quote}
My father does not smoke at home.
\end{quote}

Present: \ar{يدخن}, \emph{yudakhkhin}.

The noun \ar{تدخين} is smoking.

\begin{cousinweb}
\textbf{\ar{دخن}} (\emph{dakhkhana}) is built on the native Arabic root \emph{d-kh-n}, \textquotedblleft{}smoke\textquotedblright{} (Lane). Midkhana, a chimney, which you know, is from the same root.
\end{cousinweb}

\begin{grammarlens}[title={What you've built}]
One more word for describing what you lived and felt.
\end{grammarlens}

\subsection*{Your turn}
\begin{itemize}
\raggedright
  \item \emph{Say it:} \emph{dakhkhana}
  \item \emph{Say it:} \emph{dakhkhana}, once more
  \item \emph{Say it:} \emph{lā yudakhkhinu abī fī al-bayt.} — My father does not smoke at home.
  \item \emph{Recall:} say the Arabic for to have a cough, then the Arabic for a sprain, then say \emph{dakhkhana} again
\end{itemize}

\begin{tcolorbox}[breakable,colback=teal!4,colframe=teal!35!black,title={Before you move on}]
What does \ar{دخن} mean? (\textquotedblleft{}To smoke tobacco\textquotedblright{}.) Say it once more.
\end{tcolorbox}
\section[\texorpdfstring{\ar{اشتهى} (ishtahā)}{اشتهى (ishtahā)}]{\ar{اشتهى} (ishtahā) — to crave}

\label{lesson:AR-C517-ishtaha}



\begin{quote}
Before the new one: say the Arabic for a sprain, then the Arabic for to smoke tobacco.
\end{quote}

\subsection*{You'll want to know: \ar{اشتهى}}
\textbf{\ar{اشتهى}} — \emph{ishtahā} — \textquotedblleft{}to crave\textquotedblright{}.

\begin{quote}
\ar{اشتهت} \ar{أختي} \ar{الشوكولاتة} \ar{في} \ar{الليل}.
\end{quote}

\begin{quote}
\emph{ishtahat ukhtī ash-shūkūlāta fī al-layl.}
\end{quote}

\begin{quote}
My sister craved chocolate in the night.
\end{quote}

Present: \ar{يشتهي}, \emph{yashtahī}.

\begin{cousinweb}
\textbf{\ar{اشتهى}} (\emph{ishtahā}) is built on the native Arabic root \emph{sh-h-w}, \textquotedblleft{}to desire, to long for\textquotedblright{} (Lane). Form VIII ishtahā is to long for a thing, often food.
\end{cousinweb}

\begin{grammarlens}[title={What you've built}]
One more word for describing what you lived and felt.
\end{grammarlens}

\subsection*{Your turn}
\begin{itemize}
\raggedright
  \item \emph{Say it:} \emph{ishtahā}
  \item \emph{Say it:} \emph{ishtahā}, once more
  \item \emph{Say it:} \emph{ishtahat ukhtī ash-shūkūlāta fī al-layl.} — My sister craved chocolate in the night.
  \item \emph{Recall:} say the Arabic for a sprain, then the Arabic for to smoke tobacco, then say \emph{ishtahā} again
\end{itemize}

\begin{tcolorbox}[breakable,colback=teal!4,colframe=teal!35!black,title={Before you move on}]
What does \ar{اشتهى} mean? (\textquotedblleft{}To crave\textquotedblright{}.) Say it once more.
\end{tcolorbox}
\section[\texorpdfstring{\ar{أدمن} (admana)}{أدمن (admana)}]{\ar{أدمن} (admana) — to become addicted}

\label{lesson:AR-C517-admana}



\begin{quote}
Before the new one: say the Arabic for to smoke tobacco, then the Arabic for to crave.
\end{quote}

\subsection*{You'll want to know: \ar{أدمن}}
\textbf{\ar{أدمن}} — \emph{admana} — \textquotedblleft{}to become addicted\textquotedblright{}.

\begin{quote}
\ar{أدمن} \ar{الشاب} \ar{ألعاب} \ar{الفيديو}.
\end{quote}

\begin{quote}
\emph{admana ash-shābb alʿāb al-fīdiyū.}
\end{quote}

\begin{quote}
The young man became addicted to video games.
\end{quote}

Present: \ar{يدمن}, \emph{yudmin}.

It takes a direct object, or \ar{على}.

\begin{cousinweb}
\textbf{\ar{أدمن}} (\emph{admana}) is built on the native Arabic root \emph{d-m-n}, \textquotedblleft{}to keep to a thing\textquotedblright{} (Lane). Form IV admana is to keep at a thing without stopping.
\end{cousinweb}

\begin{grammarlens}[title={What you've built}]
One more word for describing what you lived and felt.
\end{grammarlens}

\subsection*{Your turn}
\begin{itemize}
\raggedright
  \item \emph{Say it:} \emph{admana}
  \item \emph{Say it:} \emph{admana}, once more
  \item \emph{Say it:} \emph{admana ash-shābb alʿāb al-fīdiyū.} — The young man became addicted to video games.
  \item \emph{Recall:} say the Arabic for to smoke tobacco, then the Arabic for to crave, then say \emph{admana} again
\end{itemize}

\begin{tcolorbox}[breakable,colback=teal!4,colframe=teal!35!black,title={Before you move on}]
What does \ar{أدمن} mean? (\textquotedblleft{}To become addicted\textquotedblright{}.) Say it once more.
\end{tcolorbox}
\section[\texorpdfstring{\ar{جاع} (jāʿa)}{جاع (jāʿa)}]{\ar{جاع} (jāʿa) — to get hungry}

\label{lesson:AR-C517-jaa}



\begin{quote}
Before the new one: say the Arabic for to crave, then the Arabic for to become addicted.
\end{quote}

\subsection*{You'll want to know: \ar{جاع}}
\textbf{\ar{جاع}} — \emph{jāʿa} — \textquotedblleft{}to get hungry\textquotedblright{}.

\begin{quote}
\ar{جاع} \ar{الأطفال} \ar{في} \ar{الطريق}.
\end{quote}

\begin{quote}
\emph{jāʿa al-aṭfāl fī aṭ-ṭarīq.}
\end{quote}

\begin{quote}
The children got hungry on the way.
\end{quote}

Present: \ar{يجوع}, \emph{yajūʿ}.

\begin{cousinweb}
\textbf{\ar{جاع}} (\emph{jāʿa}) is built on the native Arabic root \emph{j-w-ʿ}, \textquotedblleft{}hunger\textquotedblright{} (Lane). Jūʿ, hunger, which you know, is from the same root.
\end{cousinweb}

\begin{grammarlens}[title={What you've built}]
One more word for describing what you lived and felt.
\end{grammarlens}

\subsection*{Your turn}
\begin{itemize}
\raggedright
  \item \emph{Say it:} \emph{jāʿa}
  \item \emph{Say it:} \emph{jāʿa}, once more
  \item \emph{Say it:} \emph{jāʿa al-aṭfāl fī aṭ-ṭarīq.} — The children got hungry on the way.
  \item \emph{Recall:} say the Arabic for to crave, then the Arabic for to become addicted, then say \emph{jāʿa} again
\end{itemize}

\begin{tcolorbox}[breakable,colback=teal!4,colframe=teal!35!black,title={Before you move on}]
What does \ar{جاع} mean? (\textquotedblleft{}To get hungry\textquotedblright{}.) Say it once more.
\end{tcolorbox}
\section[\texorpdfstring{\ar{شبع} (shabiʿa)}{شبع (shabiʿa)}]{\ar{شبع} (shabiʿa) — to be full (from eating)}

\label{lesson:AR-C517-shabia}



\begin{quote}
Before the new one: say the Arabic for to become addicted, then the Arabic for to get hungry.
\end{quote}

\subsection*{You'll want to know: \ar{شبع}}
\textbf{\ar{شبع}} — \emph{shabiʿa} — \textquotedblleft{}to be full (from eating)\textquotedblright{}.

\begin{quote}
\ar{أكلت} \ar{كثيرا} \ar{وشبعت}.
\end{quote}

\begin{quote}
\emph{akaltu kathīran wa-shabiʿtu.}
\end{quote}

\begin{quote}
I ate a lot and was full.
\end{quote}

Present: \ar{يشبع}, \emph{yashbaʿ}.

\begin{cousinweb}
\textbf{\ar{شبع}} (\emph{shabiʿa}) is built on the native Arabic root \emph{sh-b-ʿ}, \textquotedblleft{}to be satisfied with food\textquotedblright{} (Lane). Shabiʿa is the opposite of jāʿa, to get hungry.
\end{cousinweb}

\begin{grammarlens}[title={What you've built}]
One more word for describing what you lived and felt.
\end{grammarlens}

\subsection*{Your turn}
\begin{itemize}
\raggedright
  \item \emph{Say it:} \emph{shabiʿa}
  \item \emph{Say it:} \emph{shabiʿa}, once more
  \item \emph{Say it:} \emph{akaltu kathīran wa-shabiʿtu.} — I ate a lot and was full.
  \item \emph{Recall:} say the Arabic for to become addicted, then the Arabic for to get hungry, then say \emph{shabiʿa} again
\end{itemize}

\begin{tcolorbox}[breakable,colback=teal!4,colframe=teal!35!black,title={Before you move on}]
What does \ar{شبع} mean? (\textquotedblleft{}To be full (from eating)\textquotedblright{}.) Say it once more.
\end{tcolorbox}
